\chapter{Ejaculation Problems}

\section*{Definition}
 Ejaculatory dysfunction includes premature ejaculation (PE), involving persistent or recurrent difficulty delaying ejaculation with personal or relationship distress, and delayed ejaculation or anejaculation, involving a marked delay or inability to ejaculate despite adequate stimulation. A time threshold alone does not diagnose PE in every person.

\section*{When to See a Doctor}
See your doctor if:
\begin{itemize}
  \item Ejaculatory difficulty is persistent and causes distress or relationship conflict
  \item Sudden onset occurs with back pain, neurological symptoms, or after starting a new medication
  \item There is associated erectile dysfunction, pain with ejaculation, or blood in semen
\end{itemize}

\section*{How It Develops}
Ejaculation is controlled by a spinal reflex (T10--L2 sympathetic, S2--4 parasympathetic and somatic) modulated by central serotonergic and dopaminergic pathways. PE is associated with reduced serotonergic tone and increased penile hypersensitivity; DE involves inadequate central excitatory signaling or excessive serotonergic inhibition.

\section*{Causes}
\begin{itemize}
  \item \textbf{Premature ejaculation:}
    \begin{itemize}
      \item Primary (lifelong): Genetic predisposition, altered 5-HT receptor sensitivity
      \item Secondary (acquired): Erectile dysfunction, prostatitis, hyperthyroidism, anxiety, relationship stress
    \end{itemize}
  \item \textbf{Delayed ejaculation:}
    \begin{itemize}
      \item Psychological: Anxiety, strict religious upbringing, masturbation style mismatch, relationship conflict
      \item Neurological: Spinal cord injury, multiple sclerosis, Parkinson disease, diabetic neuropathy
      \item Medication-induced: SSRIs (most common), antipsychotics, opioids, alpha-blockers, finasteride
      \item Endocrine: Low testosterone, hypothyroidism, prolactin disorders
    \end{itemize}
\end{itemize}

\section*{Related Symptoms}
\begin{itemize}
  \item Erectile dysfunction, anorgasmia, reduced libido, penile pain, hematospermia, infertility
\end{itemize}
\textbf{Important:} Premature ejaculation is the most common male sexual dysfunction, affecting 20--30\% of men at some point.

\section*{Medical Evaluation}
\begin{itemize}
  \item Sexual history assessing onset, frequency, context, and distress
  \item Physical examination guided by the history, including genital or neurologic examination when indicated; a digital rectal examination is not routine for every patient
  \item Psychological assessment for anxiety, depression, and relationship issues
  \item Targeted laboratory testing only when the history or examination suggests an endocrine, metabolic, or other medical contributor
\end{itemize}

\section*{Treatment Options}
\begin{itemize}
  \item \textbf{PE:} Options depend on local approval and individual factors and may include behavioral or psychological therapy, a topical anesthetic, or a prescribed SSRI; review interactions and adverse effects with a clinician
  \item \textbf{DE:} Review potentially causative medicines with a clinician and treat identified contributors; specialist-directed medicines are not routinely effective
  \item Psychological therapy for performance anxiety and relationship issues
  \item Vacuum therapy or penile vibratory stimulation for DE related to spinal cord injury
\end{itemize}

\section*{Self-Care and Home Management}
\begin{itemize}
  \item For PE: use the stop-start technique (squeeze at the glans to delay ejaculation)
  \item For DE: reduce masturbation frequency and use more stimulating techniques
  \item Communication with partner to reduce performance anxiety
  \item Review medications with a clinician (especially SSRIs, antipsychotics)
\end{itemize}

\section*{References}
\begin{thebibliography}{99}
\bibitem{ejaculation_problems:ref1} Althof SE, McMahon CG, Waldinger MD, et al. ``An Update of the International Society of Sexual Medicine's Guidelines for the Diagnosis and Treatment of Premature Ejaculation.'' J Sex Med. 2014;11(6):1392--1422.
\bibitem{ejaculation_problems:ref2} Rowland DL, Cooper SE, Schneider M. ``Premature Ejaculation: A Review.'' Int J Impot Res. 2001;13(Suppl 5):S43--S49.
\bibitem{ejaculation_problems:ref3} Sansone A, Jannini EA. ``Delayed Ejaculation and Anejaculation.'' In: IsHak WW, ed. The Textbook of Clinical Sexual Medicine. Springer; 2017.
\bibitem{ejaculation_problems:ref4} Serefoglu EC, Hellstrom WJG. ``Premature Ejaculation: Current Concepts and Treatment Options.'' Drugs. 2008;68(14):1941--1949.
\bibitem{ejaculation_problems:ref5} Waldinger MD, Quinn P, Dilleen M, et al. ``A Multinational Population Survey of Intravaginal Ejaculation Latency Time.'' J Sex Med. 2005;2(4):492--497.
\end{thebibliography}
